\SetPractica{Síntesis de aspirina}{Síntesis de aspirina}

\section{Introducción teórica}



\section{Objetivos}

\begin{itemize}
    \item Llevar a cabo la síntesis, purificación y caracterización del ácido acetilsalicílico a partir del ácido salicílico y el anhídrido acético.
\end{itemize}

\section{Materiales, equipos y reactivos}

\begin{table}[H]
\centering{\small\setstretch{1.25}
\begin{tabular}{|m{0.3\linewidth}|m{0.3\linewidth}|m{0.3\linewidth}|} \hline
    
    \thead{Equipos} & \thead{Materiales} & \thead{Reactivos} \\ \hline
    
    {Balanza \par}
    {Bomba de succión}
    
    &
    
    {Vasos de precipitado \par}
    {Probeta \par}
    {Pipetas Pasteur \par}
    {Tubos de ensayo \par}
    {Matraz kitazato \par}
    {Agitador de vidrio \par}
    {Embudo buchner \par}
    {Espátula \par}
    {Tiras de pH \par}
    {Papel filtro \par}
    {Mangueras}
    
    & 
    
    {Ácido salicílico \par}
    {Hidróxido de sodio (\ce{NaOH}) \par}
    {Hidróxido de potasio (\ce{KOH}) \par}
    {Anhídrido acético (\ce{(CH3CO)2O}) \par}
    {Acetato de etilo (\ce{CH3COOC2H5}) \par}
    {Carbonato de sodio (\ce{Na2CO3}) \par}
    {Hexano \par}
    {Ácido clorhídrico (\ce{HCl}) \par}
    {Hielo}
    
    
    \\ \hline
\end{tabular}}\end{table}

\section{Procedimientos}

\subsection{Obtención de un extracto crudo de ácido acetilsalicílico}
\begin{enumerate}
    \item En un vaso de precipitado de 100mL coloque:
    \begin{enumerate}
        \item 0.560g de ácido salicílico
        \item 1.2mL de anhídrido acético.
    \end{enumerate}
    
    \item Mezclar con una varilla los dos reactivos hasta obtener una mezcla homogénea
    \item Adicione 2.4g de carbonato de sodio a la mezcla
    \item Agitarla con una varilla de vidrio por 10 minutos.
    \item Adicione lentamente 6mL agua destilada
    \item Añade también una solución de ácido clorhídrico al 50\% hasta que el pH de la solución sea de 3.
    \item Enfríe la mezcla en un baño de hielo hasta dejar de observar cristalización.
    \begin{enumerate}
        \item El producto crudo se aísla por medio de una filtración al vacío.
    \end{enumerate}
    
    \item Lavar los cristales de ácido acetilsalicílico con agua fría.
    \begin{enumerate}
        \item Es importante que el agua esté bien fría, ya que de lo contrario el ácido acetilsalicílico se redisuelve en agua tibia.
    \end{enumerate}
    
    \item Continuar el paso de aire por succión a través de los cristales para eliminar la mayor cantidad posible de agua y la muestra de Aspirina se extiende sobre un papel de filtro para secarla al aire.
    \item Cuando el producto esté seco, pesar el producto, aunque puede contener algo de ácido sin reaccionar
    \item Calcular el rendimiento bruto.
    \begin{enumerate}
        \item Guarde al menos 0.1g de cristales de ácido acetilsalicílico para el ensayo cualitativo.
    \end{enumerate}
    
\end{enumerate}

\subsection{Purificación del ácido acetilsalicílico}
\begin{enumerate}
    \item Disolver los cristales del extracto crudo de ácido acetilsalicílico con 6mL de una mezcla caliente de acetato de etilo:hexano (60:40).
    \item Evaporar hasta aproximadamente el 70\% del volumen original
    \item Hacer esto en campana de extracción de humos
    \item Enfriar primero a temperatura ambiente y después en un baño de hielo para inducir de nuevo la formación de cristales.
    \item Filtrar los cristales por succión en un papel filtro previamente pesado
    \item Dejar secar completamente.
\end{enumerate}

\subsection{Ensayos cualitativos}
\begin{enumerate}
    \item Disolver en tres tubos de ensayo que contengan 5ml de agua, 0.1g de:
    \begin{enumerate}
        \item Ácido salicílico
        \item Producto bruto
        \item Producto purificado de ácido acetilsalicílico.
    \end{enumerate}
    
    \item Añadir una o dos gotas de solución de tricloruro de hierro al 1\% a cada uno de ellos
    \item Observar y explicar la coloración que toman.
\end{enumerate}